%%%%%%%%%%%%%%%%%%%%%%%%%%%%%%%%%%%%%%%
% Cover Letter - Pratt & Whitney
% Summer 2027 Embedded Software Engineering Intern, East Hartford CT
% Plain XeLaTeX — matches simple resume style (Spectral font)
%%%%%%%%%%%%%%%%%%%%%%%%%%%%%%%%%%%%%%%

\documentclass[11pt]{article}

\usepackage[margin=0.75in, top=0.7in, bottom=0.7in]{geometry}
\usepackage[hidelinks]{hyperref}
\usepackage{parskip}
\usepackage{fontspec}

\setmainfont{Spectral}[
  Path = ../../cv/fonts/,
  UprightFont    = Spectral-Regular.ttf,
  BoldFont       = Spectral-Bold.ttf,
  ItalicFont     = Spectral-Italic.ttf,
  BoldItalicFont = Spectral-BoldItalic.ttf
]

\setlength{\parindent}{0pt}
\setlength{\parskip}{14pt}
\linespread{1.15}
\raggedright

% No hyphenation
\hyphenpenalty=10000
\exhyphenpenalty=10000

\pagestyle{empty}
\urlstyle{same}

\begin{document}

% ── DATE ─────────────────────────────────────────────────────────────────────
\begin{flushright}
  \today
\end{flushright}

% ── BODY ─────────────────────────────────────────────────────────────────────
Dear Pratt \& Whitney Hiring Team,

I'm applying for the Summer 2027 Embedded Software Engineering Internship
on the Real-Time Embedded Systems team. I'm a Computer Science student at
RIT in the combined bachelor's
and master's program, with an expected graduation date of December 2027,
and have a strong interest in firmware and embedded software. At GE
Vernova, I built a full I/O subsystem from scratch in C on an STM32F7 for
industrial radios. On the RIT Electric Vehicle Team, I'm the vehicle
control unit (VCU) Software Lead, building real-time ThreadX firmware for
our track motorcycle, including board-level fault handling and health
monitoring across the vehicle's threads and CAN buses. In my free time I'm
building an autonomous fixed-wing flight controller on an ESP32S3, with PID
stabilization, sensor fusion, and GPS waypoint following validated across
an 80-test SITL regression suite.

What got me interested in Pratt \& Whitney specifically is that FADEC is a
real-time control loop running directly between engine sensors and the
actuators that respond to them, fuel flow, variable geometry, with
essentially no room for delay. That's the kind of sensor-to-actuator
control I like building, just at a bigger scale than a motorcycle or a
model airplane. I build that same kind of loop into the VCU, reading
sensor data off the CAN bus and driving the motor controller and
contactors in real time.

Real-time embedded C is where I spend most of my time. I've structured the
VCU's ThreadX design with a timer-triggered control thread, a health
monitoring thread, and interrupt-fed CAN queues protected by mutexes and
event flags, and I debug at the hardware level regularly with a Saleae
logic analyzer, an oscilloscope, and a J-Link Segger. I've built the kind
of hardware-abstracting software layer this team works on at a smaller
scale too, across SPI, I2C, UART, ADC, and PWM,
building both the GE Vernova I/O subsystem and EVT's hardware-agnostic
peripheral library, writing drivers directly from microcontroller
datasheets and reference manuals. At Lockheed Martin, I
also built a code-review automation tool that became a mandatory part of
the team's merge request process, so I'm comfortable with the
configuration-management and Agile-process side of embedded development
too.

I would welcome the opportunity to talk more about where I could
contribute. Thank you for your consideration.

Thanks,

Travis Brown

\end{document}
